\chapter{Introduction}

\section{Problem Statement}

Matching a customer's skin to the right foundation shade has plagued the beauty industry for decades. Conventional approaches lean almost entirely on a consultant's eye — or worse, the buyer's own guesswork — and the results are predictably uneven:

\begin{itemize}
    \item \textbf{Inconsistent Matching:} A store's fluorescent ceiling panels, a window's afternoon glare, even a consultant's personal color bias — any of these can push a recommendation off. The same customer may walk out with a different shade on two different visits.
    
    \item \textbf{Limited Accessibility:} Shade matching has traditionally demanded a physical trip to a counter staffed by a trained associate. That model shuts out anyone shopping online, living outside metropolitan areas, or simply unwilling to endure the ritual.
    
    \item \textbf{Inadequate Representation:} Brands have historically skewed their ranges toward lighter complexions. Deeper skin tones — particularly those beyond the mid-range of the Fitzpatrick scale — are routinely underserved, leaving large customer segments with few or no viable options.
    
    \item \textbf{Time and Cost:} Finding the right shade often becomes an expensive guessing game. Customers buy, test, return, and repeat — burning money and patience across multiple purchases before landing on anything close to a match.
    
    \item \textbf{Lack of Standardization:} Every brand invents its own naming scheme. "Ivory" at one label bears no guaranteed resemblance to "Ivory" at another, so cross-brand comparison becomes a minefield.
    
    \item \textbf{Subjective Undertone Detection:} Nailing the depth is only half the equation. Undertone — warm, cool, or neutral — is notoriously hard to judge by eye, and a misstep here produces that unmistakable "mask-like" mismatch.
\end{itemize}

The shift to online commerce has sharpened every one of these pain points. Shoppers ordering from a screen cannot swatch a product on their jawline; they rely on photos, reviews, and hope. Return rates in online cosmetics remain stubbornly high as a direct consequence.

\vspace{0.5cm}

Classification itself adds another layer of difficulty. The Fitzpatrick scale — originally designed to predict UV response, not cosmetic matching — compresses all of human pigmentation into six bins. That coarseness perpetuates blind spots: brands calibrated to such a narrow taxonomy naturally under-serve darker complexions.

\vspace{0.5cm}

Several AI-driven prototypes have surfaced in recent literature, yet each carries non-trivial trade-offs:

\begin{itemize}
    \item Steep computational overhead that mandates dedicated GPU infrastructure
    \item Dependence on massive, manually annotated training corpora
    \item Accuracy that degrades noticeably on darker skin tones
    \item Research-only codebases with no deployable front-end
    \item Shade databases too shallow to span the product lines customers actually encounter
\end{itemize}

What remains missing is a lightweight, inclusive pipeline — one that can ingest a single photograph, extract perceptually grounded color features, and return reliable shade recommendations without demanding a GPU cluster or a labeled dataset.

\section{Motivation}

ShadeAI grew out of a convergence: maturing open-source vision toolkits, a long-overdue industry reckoning with shade inclusivity, and the explosion of online beauty retail. Each factor reinforced the others.

\subsection{Technological Advancement}

Open-source ecosystems — OpenCV for image manipulation, scikit-learn for lightweight classification, Streamlit and FastAPI for rapid front-end and API scaffolding — have lowered the engineering barrier dramatically. Equally important, colorimetry has matured. CIEDE2000, the current gold-standard color-difference formula, models human perceptual non-uniformities with enough fidelity to anchor a shade-matching engine on solid psychophysical ground.

\subsection{Social Equity and Inclusion}

For years the cosmetics market treated darker complexions as an afterthought — fewer SKUs, narrower undertone coverage, minimal R\&D attention. Dr.\ Ellis Monk's ten-class Monk Skin Tone (MST) scale was designed explicitly to correct that imbalance. Building a recommendation system around MST does more than improve accuracy; it encodes equity into the algorithm's own classification backbone.

\subsection{E-commerce Growth}

Online cosmetics revenue has surged — a trajectory the pandemic only steepened. But digital storefronts strip away the one feedback loop shoppers historically relied on: swatching the product on their own skin. A vision-based recommendation layer can close that gap, turning a phone-camera selfie into a reliable proxy for an in-store consultation.

\subsection{Industry Demand}

Beauty conglomerates are already pouring capital into virtual try-on and AR mirror features. What many still lack, however, is a reliable perception back-end — the colorimetric engine that decides which shade to overlay. A modular, API-ready detection system slots directly into that gap.

\subsection{Research Gap}

Published literature demonstrates that high-accuracy skin classification is achievable, yet almost none of those prototypes ship with a usable interface, a real shade database, or sub-second inference on commodity hardware. The gap between a Jupyter notebook and a production endpoint is wide — and largely unaddressed.

\subsection{Personal Experience}

On a more personal note, every member of the project team has watched friends or relatives cycle through half a dozen foundations before finding anything wearable. That shared frustration — tangible, recurring, and unnecessary — anchored the project in a problem none of us had to imagine.

\subsection{Academic and Professional Development}

Academically, the project sits at the intersection of perception science, classical vision algorithms, full-stack engineering, and UX design. Few capstone scopes exercise that many disciplines simultaneously while also yielding a deployable product.

\section{Project Objectives}

ShadeAI targets a clear set of measurable goals, split into primary deliverables and supporting objectives.

\subsection{Primary Objectives}

\begin{enumerate}
    \item \textbf{Develop an Accurate Skin Tone Detection System:}
    \begin{itemize}
        \item Build a resilient image-analysis pipeline capable of handling varied lighting and camera conditions
        \item Hold recommended-shade color difference ($\Delta E$) below 3.0 — the perceptual threshold for a "good" match
        \item Classify skin against the full 10-class Monk Skin Tone scale
        \item Identify undertone (warm, cool, neutral) with heuristic-based confidence scoring
    \end{itemize}
    
    \item \textbf{Create a Comprehensive Foundation Database:}
    \begin{itemize}
        \item Aggregate shade catalogues from at least a dozen major cosmetic labels
        \item Store each shade's CIELAB coordinates for perceptually grounded matching
        \item Guarantee coverage across every MST class, especially the historically under-represented darker bins
        \item Attach structured metadata — undertone tags, MST range compatibility, product URLs
    \end{itemize}
    
    \item \textbf{Implement Production-Ready Interfaces:}
    \begin{itemize}
        \item CLI entry point for scripted batch runs and automated testing
        \item Streamlit-based web front-end for interactive, guided usage
        \item FastAPI REST endpoint returning JSON payloads for third-party integration
        \item Built-in photo-quality gate that flags blur, under-exposure, or color-cast before analysis proceeds
    \end{itemize}
    
    \item \textbf{Achieve Research-Grade Performance:}
    \begin{itemize}
        \item Meet or surpass published $\Delta E$ benchmarks (target: below 2.7)
        \item Complete end-to-end inference in under one second on a mid-range laptop CPU
        \item Operate in a zero-shot fashion — no training corpus, no fine-tuning
        \item Validate robustness across the full MST spectrum, not just lighter tones
    \end{itemize}
\end{enumerate}

\subsection{Secondary Objectives}

\begin{enumerate}
    \item \textbf{Ensure Accessibility and Usability:}
    \begin{itemize}
        \item Craft interfaces legible to users with zero technical background
        \item Ship embedded photo-capture guidelines so users know what "good input" looks like
        \item Surface actionable quality warnings (e.g., "image too dark — retake near a window")
        \item Package the system for local install, cloud deployment, and containerized (Docker) orchestration
    \end{itemize}
    
    \item \textbf{Promote Inclusivity:}
    \begin{itemize}
        \item Anchor all classification on the MST scale — no fallback to Fitzpatrick
        \item Audit the shade database to confirm deep-tone coverage is proportional, not token
        \item Stress-test the pipeline on MST-07 through MST-10 inputs specifically
        \item Ship sample images spanning the full pigmentation range so testers cannot default to light-skin-only benchmarks
    \end{itemize}
    
    \item \textbf{Enable Future Enhancements:}
    \begin{itemize}
        \item Structure the codebase into swappable modules (detection, segmentation, matching) behind clean interfaces
        \item Maintain docstrings and inline commentary at a level where a new contributor can onboard within a day
        \item Expose hooks for a future CNN-based segmentation model to replace the rule-based path
        \item Keep the shade database in a flat JSON schema that non-developers can extend
    \end{itemize}
    
    \item \textbf{Demonstrate Best Practices:}
    \begin{itemize}
        \item Adhere to SOLID design principles and PEP-8 formatting conventions
        \item Cover critical paths with unit and integration tests
        \item Deliver a self-contained documentation site alongside the source
        \item Track all changes in Git with descriptive, atomic commits
    \end{itemize}
\end{enumerate}

\section{Scope of the Project}

\subsection{In Scope}

The current release of ShadeAI targets the feature set below. Each capability is fully implemented, tested, and documented.

\begin{enumerate}
    \item \textbf{Image Processing Pipeline:}
    \begin{itemize}
        \item White-balance normalization via the Gray World assumption
        \item On-the-fly conversion across RGB, CIELAB, HSV, and YCbCr spaces
        \item Six-metric quality gate (brightness, blur, color cast, face presence, resolution, exposure)
        \item JPEG and PNG ingestion out of the box
    \end{itemize}
    
    \item \textbf{Face Detection:}
    \begin{itemize}
        \item Primary detector: OpenCV Haar Cascade (frontal face)
        \item Graceful fallback path when the primary cascade yields no hits
        \item Bounding-box extraction followed by geometric sanity checks
    \end{itemize}
    
    \item \textbf{Skin Segmentation:}
    \begin{itemize}
        \item Dual-space thresholding (HSV + YCbCr) for robust skin-pixel isolation
        \item Morphological opening and closing to suppress speckle noise
        \item Pixel-count validation to reject frames with insufficient skin area
    \end{itemize}
    
    \item \textbf{Feature Extraction:}
    \begin{itemize}
        \item ITA° computation from mean $L^*$ and $b^*$ of the segmented region
        \item First-order CIELAB statistics: mean, standard deviation, and range of $L^*$, $a^*$, $b^*$
        \item Histogram-level distribution profiling for outlier detection
    \end{itemize}
    
    \item \textbf{Classification:}
    \begin{itemize}
        \item Ten-bin MST assignment driven by calibrated ITA° thresholds
        \item Undertone inference via $a^*$/$b^*$ ratio heuristics
        \item Per-prediction confidence score reflecting feature-space margin to the nearest class boundary
    \end{itemize}
    
    \item \textbf{Shade Recommendation:}
    \begin{itemize}
        \item CIEDE2000-based ranking of every shade in the database against the user's extracted CIELAB vector
        \item Top-five results returned, sorted by ascending $\Delta E_{00}$
        \item Optional undertone filter to suppress cross-undertone suggestions
        \item Each result annotated with brand, product line, shade name, and purchase link where available
    \end{itemize}
    
    \item \textbf{Foundation Database:}
    \begin{itemize}
        \item 187 shades spanning 14 brands — from drugstore to prestige tier
        \item Lab-measured CIELAB triplets for each entry
        \item Undertone tag and MST-range compatibility baked into every record
        \item Flat JSON schema; new shades can be appended without touching application code
    \end{itemize}
    
    \item \textbf{User Interfaces:}
    \begin{itemize}
        \item CLI that dumps full analysis (MST class, CIELAB stats, top-five shades) to stdout
        \item Streamlit app with drag-and-drop upload, live preview, and expandable result cards
        \item FastAPI endpoint returning structured JSON — ready for downstream services
        \item Integrated quality checker that runs before analysis and warns on suboptimal input
    \end{itemize}
    
    \item \textbf{Documentation:}
    \begin{itemize}
        \item README covering installation, quickstart, and configuration
        \item API reference with curl-ready examples for every endpoint
        \item Inline docstrings following NumPy-style conventions
        \item Reproducible accuracy report comparing results against published benchmarks
    \end{itemize}
\end{enumerate}

\subsection{Out of Scope}

The items below fall outside the present release. Several are active candidates for a follow-up iteration:

\begin{enumerate}
    \item Deep learning models for segmentation and classification
    \item Real-time webcam capture and analysis
    \item Virtual try-on with augmented reality
    \item Mobile applications (iOS, Android)
    \item Multi-face detection and analysis
    \item Skin condition analysis (acne, wrinkles, etc.)
    \item Personalized skincare recommendations
    \item Integration with e-commerce platforms
    \item User accounts and history tracking
    \item Social media sharing features
\end{enumerate}

\subsection{Limitations}

Every vision-based system carries constraints. The ones relevant here are:

\begin{enumerate}
    \item Input images must be reasonably well-lit; strong shadows or color casts degrade accuracy
    \item The Haar detector is tuned for near-frontal poses — extreme angles will fail silently
    \item Heavy cosmetic coverage or Instagram-style filters can shift the extracted CIELAB centroid significantly
    \item The shade catalogue, while diverse at 187 entries, does not yet cover niche or regional brands
    \item Undertone classification relies on hand-crafted $a^*/b^*$ thresholds rather than a learned model
    \item Video streams are not supported; the pipeline expects single-frame inputs
    \item Multi-face scenes are unsupported — only the largest detected face is analyzed
\end{enumerate}

\section{Organization of the Report}

The report follows a six-chapter arc, bookended by two appendices.

\textbf{Chapter 1: Introduction} frames the problem landscape, articulates why an AI-driven shade-matching tool is overdue, and pins down the project's objectives and boundaries.

\textbf{Chapter 2: Literature Review} surveys prior work in skin-tone classification, perceptual colorimetry, and cosmetic recommendation engines, surfacing the specific shortcomings that ShadeAI is designed to address.

\textbf{Chapter 3: System Design and Methodology} walks through the six-stage pipeline architecture, detailing each module's algorithm, data flow, and the engineering trade-offs that shaped it.

\textbf{Chapter 4: Implementation and Results} covers the development environment, deployment stack, and experimental evaluation — including per-MST-class accuracy breakdowns and head-to-head comparisons with published baselines.

\textbf{Chapter 5: Conclusion and Future Scope} distils the project's contributions, acknowledges its boundaries, and charts a roadmap for subsequent iterations.

\textbf{Chapter 6: References} compiles every cited paper, standard, software library, and web resource.

\textbf{Appendix A: List of Papers} records publications, submissions, or communications arising from this project.

\textbf{Appendix B: Plagiarism Report} reproduces the mandatory plagiarism and AI-content screening certificates.

\vspace{0.5cm}

Taken together, the chapters trace ShadeAI from problem definition through algorithmic design, implementation, quantitative evaluation, and forward-looking research directions — illustrating how classical vision techniques, modern color science, and thoughtful software engineering can converge on a socially meaningful, commercially viable product.
